\subsection{Concept Two: EVA}
\begin{center}
    \textit{Short for Evangelion, draws from Roman times and means 'messenger.' Reflects the satellite's mission to carry knowledge and insights of the cosmos back to Earth.}
\end{center}
\subsubsection{Abstract} 
The second concept - EVA, is a Highly Elliptical Orbit (HEO) travelling satellite that ranges from LEO to GEO attitude. EVA offers a competitive balance between serviceability and an undisturbed photographic environment. The fundamental challenge with the Hubble and James Webb Telescope arises from their distinct operational environments: Hubble's close proximity with Earth's atmosphere introduces limitation, while James's abdicant distance disallows regular service. EVA uses the perigee of it's orbit to operate communication tasks, and it's apogee for photographing our universe without the Earth, commercial satellite, or atmosphere's intervene. The satellite will use modern level equipment that emphasize on autonomous operation - Tailoring to the complex nature of long distance communication and latency. The goal of this concept is to expand the theoretical limits of our telescope by utilizing it's unique orbiting type, providing extraordinary images in a serene environment while obtaining the redundancy of servicing - Balancing photographic performance with sustainable operational redundancy at the cost of greater price and innovation. 
\subsubsection{Orbit Type}
A Highly Elliptical Orbit (HEO) offers the optimal trajectory that balances a stable deep-space photographic environment and the accessibility of servicing. Unlike the Hubble Space Telescope, which is affected by Earth’s albedo IR radiation, atmospheric drag, light pollution, and interference from commercial satellites, as well as a time constraint on constant surveying at a stationary point. (View of the celestial objects gets physically blocked as Hubble orbits around Earth). EVA satellite in HEO circumvents these limitations. At apogee, the Earth no longer obstructs its field of view, allowing a greater Field Of View (FOV) for extended observation periods and greater flexibility for continuous sky surveys. Furthermore, future low-cost servicing mission could be conducted during periods of orbital perigee, offering greater redundancy compared to James Webb while also reaping the benefits of operating at a distance far away from earth for optimal imaging.  

However, these advantages come with trade-offs, and there are reasons why HEO are not the usual trends of satellites. HEO orbit requires a much greater delta-V and time to transition \& maintain it's highly epileptical shape - Contributing to it's higher launch and operational costs (Excludes servicing cost as it is designed for servicing). The orbit trajectory also presents significant transitional challenges, such as frequent exposure to Van Allen Belt radiation, temperature fluctuations, drastic velocity changes, limited communication windows, and complex power management. The satellite as a result shall also be equipped with advanced radiation hardening, active thermal control systems to handle extreme temperature variations, and robust mechanical structures to endure the dynamic stresses of orbit transitions. Autonomous operations will be critical due to communication delays, but the ability to service the satellite makes EVA a cost-effective and sustainable solution for long-term deep-space exploration.
[N1,2,3]
\subsubsection{Camera}

\vspace{1.5mm}

EVA satellite, operating in a Highly Elliptical Orbit (HEO), requires a specialized camera system to fully exploit the unique benefits of this orbit. To address the varying environmental challenges and maximize imaging capabilities, several critical features have been integrated into the camera design.

\vspace{1.5mm}

First, the camera must handle significant variations in light conditions as the satellite moves from perigee to apogee. For this, wide dynamic range sensors are used, allowing clear imaging despite changes in brightness. Additionally, the system employs radiation-hardened CMOS sensors to endure the high radiation levels encountered in the Van Allen Belts, ensuring long-term performance. Furthermore, to maintain sharp imaging despite the varying velocities in HEO, the camera is equipped with adaptive optics and image stabilization systems, which correct for motion blur during rapid orbital transitions. Furthermore, multi-spectral imaging capabilities covering infrared (IR), visible, and ultraviolet (UV) wavelengths enable comprehensive deep-space observations.

\vspace{1.5mm}

Given the extended observation periods possible at apogee, a large aperture of 2-3 meters has been selected to improve light-gathering power, thus enhancing resolution and detail. Cryogenic cooling is also implemented to reduce thermal noise in infrared imaging, crucial for detecting faint celestial objects. Lastly, the system includes onboard autonomous data processing powered by AI, optimizing data compression and prioritization to efficiently use limited communication windows. This approach ensures that critical data is transmitted back to Earth, maximizing the satellite’s scientific output while operating independently for extended periods.


\vspace{1.5mm}

\subsubsection{Control \& Communications Features}
EVA shall utilize an autonomous decision making intelligence that controls the satellite during communication downlink time. The Eva shall utilize similar control surfaces to the NOVAE, where a HRG and FGS will be implemented. A greater SCAT thrusters are also needed to navigate and maintain the high-speed nature of it's orbital trajectory. A larger flash memory will need to be mounted onboard due to the expected communication downlink period; Data shall be store locally until connection with ground station is available.  HRG will also serves as an orientation tool to direct the facing of the telescope. Upon the perigee, the satellite shall conduct bidirectional high-speed data transfer - Transferring image data while uploading ground stations mission command concurrently. This bridges the connection between the user and the satellite, allow the user to direct and take images of stars, planets, and galaxies. [N1,N3]

\subsubsection{Sustainability}
The user has the option to terminal this operation by redirecting it's trajectory into the Earth's atmosphere over a period of time. Sustainable material and manufacturing approach shall also be prenticed (More on sustainability on concept 3). [N5]
